\documentclass[11pt,a4paper]{article}
\usepackage[utf8]{inputenc}
\usepackage{amsmath,amssymb}

\setlength{\topmargin}{-0.5in}
\setlength{\textheight}{9.2in}
\setlength{\oddsidemargin}{0in}
\setlength{\evensidemargin}{0in}
\setlength{\textwidth}{6.5in}

\newtheorem{theorem}{Theorem}
\newtheorem{lemma}{Lemma}

\title{\textbf{Derivation D2: Logarithmic Height Bound for AVL Trees}\\
\large Exact Derivation of $h < 1.4405 \log_2(n + 2) - 0.3277$ via Minimal Fibonacci Trees}
\author{\textbf{Saumya (250001064)} -- Role P2 (AVL Tree Engineer)\\
\textit{CS 253 / MA 253 Data Structures \& Algorithms Lab}}
\date{21 November 2026}


\newenvironment{proof}{\noindent\textbf{Proof:}\enspace}{\hfill$\blacksquare$\par\medskip}
\begin{document}

\maketitle

\section{Definition and Structural Invariant}
An AVL tree is a binary search tree in which for every node $x$, the heights of its left and right subtrees differ by at most $1$. Formally, let $h(x)$ denote the height of node $x$, defined as the number of edges on the longest path from $x$ to a descendant leaf:
\begin{equation}
h(\text{nullptr}) = -1, \quad h(\text{leaf}) = 0, \quad h(x) = 1 + \max(h(x.\text{left}), h(x.\text{right}))
\end{equation}
The balance factor $\text{bf}(x)$ is defined as:
\begin{equation}
\text{bf}(x) = h(x.\text{left}) - h(x.\text{right})
\end{equation}
The AVL balance invariant requires that for every node $x$ in the tree:
\begin{equation}
\text{bf}(x) \in \{-1, 0, +1\}
\end{equation}

\section{Minimal AVL Trees and Fibonacci Recurrence}
To determine the maximum height $h$ that an AVL tree containing $n$ nodes can attain, we analyze the dual problem: \textit{What is the minimum number of nodes, $N(h)$, required to construct an AVL tree of height $h$?}

For an AVL tree of height $h$ to contain the absolute minimum number of nodes, its root must have one child subtree of minimal height $h - 1$ and another child subtree of minimal height $h - 2$:
\begin{equation}
N(h) = 1 + N(h - 1) + N(h - 2) \quad \text{for } h \ge 1
\end{equation}
with boundary conditions:
\begin{equation}
N(-1) = 0 \quad (\text{empty tree}), \qquad N(0) = 1 \quad (\text{single-node tree})
\end{equation}

Computing the initial terms:
\begin{align*}
N(1) &= 1 + N(0) + N(-1) = 1 + 1 + 0 = 2 \\
N(2) &= 1 + N(1) + N(0) = 1 + 2 + 1 = 4 \\
N(3) &= 1 + N(2) + N(1) = 1 + 4 + 2 = 7 \\
N(4) &= 1 + N(3) + N(2) = 1 + 7 + 4 = 12 \\
N(5) &= 1 + N(4) + N(3) = 1 + 12 + 7 = 20
\end{align*}

\begin{lemma}
For all $h \ge -1$, $N(h) = F_{h+3} - 1$, where $F_k$ is the $k$-th Fibonacci number ($F_0=0, F_1=1, F_2=1, F_3=2, \dots$).
\end{lemma}
\begin{proof}
By mathematical induction on $h$:
\begin{itemize}
    \item \textbf{Base cases:} For $h = -1$, $N(-1) = 0$ and $F_{-1+3} - 1 = F_2 - 1 = 1 - 1 = 0$. For $h = 0$, $N(0) = 1$ and $F_{0+3} - 1 = F_3 - 1 = 2 - 1 = 1$. Both hold.
    \item \textbf{Inductive step:} Assume $N(k) = F_{k+3} - 1$ holds for all $k < h$. Then:
    \begin{align*}
    N(h) &= 1 + N(h - 1) + N(h - 2) \\
         &= 1 + (F_{h+2} - 1) + (F_{h+1} - 1) \\
         &= F_{h+2} + F_{h+1} - 1 \\
         &= F_{h+3} - 1
    \end{align*}
    by the standard Fibonacci identity $F_{k} = F_{k-1} + F_{k-2}$. Hence the identity holds for all $h \ge -1$.
\end{itemize}
\end{proof}

\section{Exact Closed-Form Height Bound}
By Binet's formula, the $k$-th Fibonacci number is given by:
\begin{equation}
F_k = \frac{\varphi^k - \psi^k}{\sqrt{5}}
\end{equation}
where $\varphi = \frac{1 + \sqrt{5}}{2} \approx 1.6180339887$ (Golden Ratio) and $\psi = \frac{1 - \sqrt{5}}{2} \approx -0.6180339887$.

Since $|\psi| < 1$, for all $k \ge 1$ we have $|\psi^k / \sqrt{5}| \le \frac{1}{\sqrt{5}} < 0.44721$. Thus:
\begin{equation}
F_{h+3} > \frac{\varphi^{h+3}}{\sqrt{5}} - 1
\end{equation}
Substituting into the minimal node relation for an AVL tree with $n$ nodes:
\begin{equation}
n \ge N(h) = F_{h+3} - 1 > \frac{\varphi^{h+3}}{\sqrt{5}} - 2
\end{equation}
Rearranging gives:
\begin{equation}
n + 2 > \frac{\varphi^{h+3}}{\sqrt{5}}
\end{equation}
Taking base-2 logarithms on both sides:
\begin{align}
\log_2(n + 2) &> (h + 3)\log_2(\varphi) - \frac{1}{2}\log_2(5) \\
h + 3 &< \frac{\log_2(n + 2) + \frac{1}{2}\log_2(5)}{\log_2(\varphi)} \\
h &< \frac{1}{\log_2(\varphi)} \log_2(n + 2) + \left[ \frac{\log_2(5)}{2\log_2(\varphi)} - 3 \right]
\end{align}

\noindent\fbox{%
\parbox{0.96\textwidth}{%
\textbf{Numerical Constant Evaluation:}\\
$\varphi \approx 1.6180339887 \implies \log_2(\varphi) \approx 0.6942419136$\\
$\frac{1}{\log_2(\varphi)} \approx 1.4404200904$\\
$\frac{\log_2(5)}{2\log_2(\varphi)} - 3 = \frac{2.3219280949}{2 \times 0.6942419136} - 3 = 1.672275 - 3 = -1.327725$\\
Rounding upward for a strictly conservative ceiling bound gives:
\begin{equation}
h < 1.4405 \log_2(n + 2) - 0.3277
\end{equation}
}%
}

\section{Implications for Operation Bounds}
For any AVL tree storing $n$ events:
\begin{itemize}
    \item For $n = 1{,}000$, $h < 1.4405 \times 9.9687 - 0.3277 \approx 14.03 \implies h \le 14$.
    \item For $n = 1{,}000{,}000$, $h < 1.4405 \times 19.9316 - 0.3277 \approx 28.38 \implies h \le 28$.
\end{itemize}
All operations ($\text{insert}$, $\text{deleteMin}$, $\text{cancel}$, and $\text{successor}(t)$) traverse at most $h + 1 \le 29$ pointers, strictly guaranteeing $O(\log n)$ worst-case time with no performance degradation.
\end{document}
